% =====================================================================
% Argos - Literature Review
% A Multi-Agent LLM Trading Research Desk
% IEEE conference format (IEEEtran), Overleaf-ready
%
% Compile: pdflatex -> bibtex -> pdflatex x 2
% =====================================================================
\documentclass[conference]{IEEEtran}

\usepackage{cite}
\usepackage{amsmath,amssymb,amsfonts}
\usepackage{graphicx}
\usepackage{booktabs}
\usepackage{array}
\usepackage{tabularx}
\usepackage{enumitem}
\usepackage{url}
\usepackage[colorlinks=true,linkcolor=black,citecolor=black,urlcolor=black]{hyperref}

% ---- Shorthand ------------------------------------------------------
\newcommand{\sys}{Argos}
\newcolumntype{L}[1]{>{\raggedright\arraybackslash}p{#1}}

\begin{document}

\title{\sys{}: A Multi-Agent Large Language Model Trading Research Desk\\
\large Literature Review (v1.0)}

\author{\IEEEauthorblockN{Kishor Timilsena}
\IEEEauthorblockA{\textit{Fusemachines AI Fellowship 2026}\\
Kathmandu, Nepal\\
kishor.timilsena@example.com}}

\maketitle

% =====================================================================
\begin{abstract}
This review collects and organizes the research behind \sys{}, a multi-agent research desk that helps a trader decide whether to buy, sell, or hold a stock. The desk uses five large language model (LLM) roles and one candlestick forecasting model called Kronos, and a Risk Manager that can veto any trade. The review starts with the finance ideas the project depends on: how markets are studied, what candlesticks and technical analysis are, and how news text moves prices. It then covers backtesting, the method used to test a trading idea on past data without cheating, and the well known ways backtests lie. Next it covers forecasting models, from the Transformer to modern time-series foundation models, and then LLMs in finance, LLM agents and teams, and the research on getting models to argue with each other. Every paper is tied back to what it gives this project. The review closes with the gap that no paper covers: a system where a forecast model's uncertainty, read by an independent Risk Manager with veto power, can overrule a confident debate. That gap is the research claim of \sys{}.
\end{abstract}

\begin{IEEEkeywords}
multi-agent systems, large language models, decision support, algorithmic trading, foundation models, time-series forecasting, backtesting, risk management, adversarial debate
\end{IEEEkeywords}

% =====================================================================
\section{Introduction}
% =====================================================================

\subsection{Why This Review Exists}

\sys{} sits at the meeting point of three research areas that do not usually talk to each other. The first is quantitative finance: how traders read price data, how they test ideas on history, and how they control risk. The second is forecasting: models that predict the next numbers in a sequence. The third is LLM agents: programs where one or more language models play roles, use tools, and work together. A useful review has to cover all three, because the project borrows something real from each.

The review is organized the same way the project is built. Sections~\ref{sec:finance} and~\ref{sec:backtest} cover the finance side: what the desk is looking at, and how we will test it honestly. Sections~\ref{sec:forecasting} and~\ref{sec:llmfinance} cover the models: what forecasts the numbers, and what the LLMs do. Sections~\ref{sec:agents} and~\ref{sec:debate} cover the agents: how single agents work, how teams work, and what happens when models argue. Section~\ref{sec:risk} covers risk. Sections~\ref{sec:gap} and~\ref{sec:mapping} put it all together: what the literature does not yet cover, and how each paper maps to a design decision in \sys{}.

\subsection{How Papers Were Chosen}

Three rules. First, a paper must either be a direct building block of the system (for example Kronos, which is literally a component), or a well known result in its area that defines what good work looks like (for example the Efficient Markets Hypothesis, or the surveys of backtest overfitting). Second, areas were covered broadly on purpose: finance classics, technical analysis, textual finance, forecasting architectures, time-series foundation models, financial LLMs, LLM agents, debate research, and decision psychology. A project that mixes these areas needs a reading list that mixes them too. Third, every paper here earns its place by doing something for the project. Sections~\ref{sec:gap} and~\ref{sec:mapping} make that contribution explicit, paper by paper.

% =====================================================================
\section{Finance Ideas the Project Depends On}
\label{sec:finance}
% =====================================================================

\subsection{Can Public Information Beat the Market?}

The starting question for any trading system is whether it should expect to find anything at all. The Efficient Markets Hypothesis (EMH) says that prices already reflect public information, so rules built on public information should not keep making money after costs \cite{fama1970efficient}. This is not a claim that nobody ever profits. It is a claim that any repeatable edge must come from information that is either not public, not yet in the price, or processed better than everyone else does.

\sys{} takes the EMH seriously as the null hypothesis. The project does not promise a profitable strategy. It builds a decision-support desk, measures it against simple baselines like buy and hold, and reports honestly when it loses. The Adaptive Markets Hypothesis (AMH) adds the missing piece: markets move through phases, and the inefficiencies available in one phase can vanish in the next \cite{lo2004adaptive}. This is why the project evaluates over one fixed window plus one stress window, and why the roadmap treats multi-regime testing as future work rather than pretending one window proves anything.

\emph{What the project takes from this:} the evaluation plan in \sys{} treats buy-and-hold as the bar that must be cleared, and treats any result as window-specific, exactly as the EMH and AMH literature says it should.

\subsection{Candles and Technical Analysis}

Price data comes in candles. One candle is one period of trading, usually one day here, and records five numbers: the open price, the highest price, the lowest price, the close price, and the volume traded. A series of candles is called OHLCV data, and it is the raw diet of every technical method.

Technical analysis is the practice of reading patterns in this data: moving averages crossing, support and resistance levels, momentum. The evidence for it is genuinely mixed, and the two most cited studies frame the debate. Brock, Lakonishok, and LeBaron found that simple moving-average and trading-range rules had real predictive power on ninety years of Dow Jones data \cite{brock1992technical}. A decade later, Park and Irwin surveyed the modern literature and found that while some studies report positive results, most suffer from data problems, testing biases, and risk adjustment issues, so the honest summary is that profitability evidence is mixed and fragile \cite{park2013intraday}.

\emph{What the project takes from this:} \sys{} deliberately does \emph{not} compute RSI or MACD and hand them to an LLM. Instead, the technical leg is a forecasting model over the raw candles. This sidesteps the pattern-reading debate entirely and keeps the desk's technical view grounded in numbers rather than in a language model's opinion about chart patterns.

\subsection{News and Text as Market Data}

Long before LLMs, researchers showed that text moves prices. Tetlock showed that the tone of Wall Street Journal columns predicts short-term market moves, with high media pessimism followed by a temporary dip and a rebound \cite{tetlock2007giving}. Antweiler and Frank studied hundreds of thousands of stock message-board posts and found that, although most talk is noise, the volume and the bullishness of messages have measurable predictive power for volatility and, more weakly, for returns \cite{antweiler2004text}. Loughran and McDonald showed why generic sentiment dictionaries fail on finance text: words like \emph{liability} or \emph{capital} are neutral in daily life but loaded in filings, and a finance-specific word list predicts stock returns far better than a general one \cite{loughran2011liability}.

LLMs changed the scale of this work. Lopez-Lira and Tang gave ChatGPT headlines and asked whether the news was good, bad, or irrelevant for a stock's next-day price, and the model's scores predicted next-day returns better than traditional sentiment methods \cite{lopezlira2023chatgpt}. This result is the single strongest justification for the Sentiment Analyst role in \sys{}: an LLM can read a headline the way Loughran and McDonald said a finance-literate reader would, without a hand-built dictionary.

\emph{What the project takes from this:} the Sentiment Analyst is a real role with a real evidence base, not decoration. The literature also explains its main risk: tone in text is easy to over-read, so the agent must work only from provided headlines and rate its confidence honestly when the news is thin.

% =====================================================================
\section{Backtesting: Testing Ideas on History Without Cheating}
\label{sec:backtest}
% =====================================================================

A backtest replays a trading idea on past data: at each historical date the system decides, then the actual outcome is revealed and scored. Done right, it is the only cheap way to know if an idea has any worth. Done wrong, it is a machine for fooling yourself.

\subsection{The Ways Backtests Lie}

Lopez de Prado's \emph{Advances in Financial Machine Learning} is the standard reference on doing this correctly, and its core warnings drive the \sys{} evaluation design \cite{lopezdeprado2018advances}. Bailey, Borwein, Lopez de Prado, and Zhu showed the mechanism behind the worst failure: if you try enough strategy variations on the same historical data, one of them will look amazing purely by chance, a problem they call backtest overfitting \cite{bailey2014backtest}. The same authors later proposed the Deflated Sharpe Ratio, a way to correct a reported performance number for how many attempts were made before it \cite{bailey2014deflated}.

\emph{What the project takes from this:} three concrete rules. The baselines and ablations in \sys{} are fixed in advance, before any results exist, which caps the number of attempts. The stress-test window is chosen independently of the main window. And the document commits to publishing failures with the same prominence as wins, because a report with only wins is the signature of an overfitted backtest.

\subsection{No Lookahead and Walk-Forward Testing}

The most basic backtest error is lookahead: letting the system see data from after the date it is deciding about. \sys{} enforces this in code rather than by discipline: one loader function is the only code allowed to read raw data, it takes an explicit as-of date, and it refuses to return anything newer. The walk-forward protocol then moves through the window day by day, deciding before each outcome is revealed \cite{lopezdeprado2018advances}.

\subsection{Scoring Forecasts and Comparing Systems}

Two measurement papers matter here. Hyndman and Koehler proposed MASE, the Mean Absolute Scaled Error, which expresses a forecast's error relative to a naive baseline, making error scores comparable across series \cite{hyndman2006another}. Kronos itself reports MASE and WQL, the Weighted Quantile Loss, so the project scores its forecasting component in the same units the forecasting literature uses. Diebold and Mariano provided the standard statistical test for asking whether one forecasting system is genuinely better than another, or just luckier on this sample \cite{diebold1995comparing}. The Sharpe ratio, proposed by Sharpe as a reward-to-variability measure, remains the standard way to compare risk-adjusted returns \cite{sharpe1994sharpe}.

\emph{What the project takes from this:} every metric in the evaluation plan has a source and a purpose: MASE and WQL for forecast quality, the Diebold-Mariano test when comparing against baselines, Sharpe and maximum drawdown for the simulated position series, and the Deflated Sharpe correction as the honest lens over all of it.

% =====================================================================
\section{Forecasting Prices with Models}
\label{sec:forecasting}
% =====================================================================

\subsection{The Transformer}

The Transformer architecture, introduced by Vaswani et al., replaced recurrence with attention and became the base for nearly all modern sequence models \cite{vaswani2017attention}. Everything in this section and the next section, including every LLM and every time-series foundation model in the project, is a Transformer under the hood.

\subsection{Time-Series Transformers}

Forecasting research adapted the Transformer to time series in two waves. In the first wave, PatchTST treated a time series like a sentence, cutting it into patches and forecasting with attention, and beat strong previous models on standard benchmarks \cite{nie2023patchtst}. In the second wave, researchers stopped training from scratch. Zhou et al. showed that a pretrained language model, with small changes, could be repurposed for time-series work and match specialized models, an idea they summarized as one model fits all \cite{zhou2023onefitsall}.

\subsection{Time-Series Foundation Models}

That repurposing idea grew into time-series foundation models: large models pretrained on huge amounts of series data, ready to forecast a new series with no training. Lag-Llama produced full probability distributions over future values, not single point predictions \cite{rasul2023lagllama}. Chronos tokenized numeric series into discrete tokens and showed that language-model machinery, trained on time series, forecasts well without seeing the domain \cite{ansari2024chronos}. Both made the same conceptual move: treat a series like a language, and sample the future instead of computing it.

\subsection{Kronos: A Foundation Model for Candles}

Kronos is the first open-source foundation model built specifically for financial candles \cite{shi2025kronos}. It was pretrained on candle data from over 45 global exchanges. A special tokenizer turns each candle into discrete tokens, and a decoder-only Transformer generates future candle tokens one after another, which decode back into future prices. Two properties matter for \sys{}. First, because decoding is stochastic, Kronos produces a \emph{distribution} of possible futures: run it many times and you get many paths, whose spread measures how sure the model is. Second, Kronos has no language interface at all. It cannot explain a forecast, and it cannot take instructions.

\emph{What the project takes from this:} Kronos is the Technicals Analyst's engine, and the spread of its sampled paths, especially the gap between the 90th and 10th percentile, becomes the Risk Manager's main numeric signal. The fact that Kronos cannot explain itself is handled by design: a separate LLM narrates its numbers faithfully and is forbidden from inventing reasons. The other foundation models matter as context and as possible alternatives, but Kronos is the one trained on exactly the object this project cares about: candles.

% =====================================================================
\section{LLMs in Finance}
\label{sec:llmfinance}
% =====================================================================

Three projects defined what a finance LLM looks like. BloombergGPT showed that mixing general text with large amounts of finance text produces a model better at finance tasks than a general model of the same size \cite{wu2023bloomberggpt}. FinGPT took the open-source path: an automatic data pipeline and lightweight fine-tuning so that financial LLMs can be updated cheaply as new data arrives \cite{yang2023fingpt}. FinRobot organized LLMs into an agent platform for financial applications, with layers for data, algorithms, and application-specific agents \cite{yang2024finrobot}. MarketSenseAI built a queryable service that combines news, fundamentals, and macro data through LLMs to explain stock selections \cite{fatouros2024marketsenseai}.

These systems share a lesson: the LLM is good at reading and combining evidence, not at doing arithmetic or knowing the future. Lopez-Lira and Tang's result \cite{lopezlira2023chatgpt} is the cleanest evidence of real predictive value in the reading role, and it came precisely from a narrow task: score the news, do not build a whole strategy around the model's vibes.

\emph{What the project takes from this:} in \sys{}, every LLM role gets a narrow job. Analysts read provided data and report signals. Advocates argue one side. The Portfolio Manager weighs the argument. None of them is asked to compute indicators or forecast prices, jobs that belong to code and to Kronos.

A further warning comes from the hallucination literature: language models generate fluent but false content, and surveys of the problem show it is structural, not a bug that prompting fully removes \cite{ji2023hallucination}. \sys{} answers with structure: outputs must pass a fixed schema, every number in a rationale must trace to provided inputs, and sampled outputs are audited by hand during development.

% =====================================================================
\section{LLM Agents and Agent Teams}
\label{sec:agents}
% =====================================================================

\subsection{From a Single Agent to a Team}

The base pattern for an LLM agent is ReAct: the model interleaves reasoning steps with actions like tool calls, observing each result before the next step \cite{yao2022react}. Park et al. showed that a population of LLM agents with memory, routines, and social interaction produces surprisingly believable behavior, establishing that role and memory design, not just model quality, shapes what agents do \cite{park2023generative}. ChatDev showed that a team of role-playing agents, given a communication protocol, can complete real software tasks end to end \cite{qian2023chatdev}. Guo et al. survey this field and identify the design choices that actually matter: how agents are assigned roles, how they communicate, and how their work is stopped and checked \cite{guo2024llmmultiagents}.

\emph{What the project takes from this:} \sys{}'s agents are deliberately simple. Each one is a single structured LLM call with a role prompt and a fixed input bundle. The project uses none of the autonomous tool-using machinery, because the desk's workflow is fixed and auditable by design. What it borrows is the deeper finding from this literature: role separation and communication structure are first-order design choices that change output quality.

\subsection{LLM Trading Agents}

FinMem gives a single LLM agent a layered memory and a fixed trader character, and shows improved trading results against baselines \cite{yu2023finmem}. FinRL represents the earlier reinforcement-learning approach to the same problem: agents trained on market data to trade directly \cite{liu2023finrl}. TradingAgents, from Tauric Research, is the closest system to \sys{}: fundamental, sentiment, technical, and news analyst agents; bull and bear researcher agents that debate; a risk team; and a trader agent that synthesizes, all organized to mirror a real trading firm \cite{xiao2024tradingagents}. Ding et al. survey LLM trading agents and find the field converging on the same pipeline shape: analysts feed a decision step, with limited genuine debate and no independent risk gate \cite{ding2024llmtrading}.

\emph{What the project takes from this:} the desk metaphor itself, validated by TradingAgents. The differences are the point: \sys{} replaces the LLM-reads-indicators technical leg with Kronos, gives the Risk Manager a structural veto instead of advisory influence, and targets decision support with audit trails rather than autonomous trading.

% =====================================================================
\section{Debate: Making Models Argue}
\label{sec:debate}
% =====================================================================

\subsection{The Debate Literature}

Du et al. showed the foundational result: when multiple LLM instances answer a question, then see each other's answers and revise over several rounds, factuality and reasoning improve on math, strategy, and factual tasks \cite{du2023debate}. Liang et al. added role prompting to the debate, giving each participant an assigned perspective, and found it improves both reasoning and the diversity of answers compared with self-consistency methods \cite{liang2023divergent}. ChatEval used multiple debating agents as judges of text quality and found they better match human judgments than a single model scoring alone \cite{chan2023chateval}.

Smit et al. then asked the question \sys{} cares about: which debate format should you actually use? They compared formats, agent counts, and round counts across tasks and found that choices matter, that more rounds do not always help, and that debate always costs more compute \cite{smit2024mad}. A separate line of work studies LLM-as-judge: using a strong model to grade other models' outputs. Zheng et al. showed such judges agree with human preferences at useful rates but also documented known biases, like preferring a model's own outputs \cite{zheng2023judging}.

\emph{What the project takes from this:} \sys{} uses the smallest debate with real teeth: Bull argues, Bear rebuts the Bull's specific claims, Portfolio Manager weighs both. The formats and round counts identified by the debate literature are the project's tuning knobs, held at their cheapest settings. The LLM-as-judge idea appears only in evaluation, as a checked instrument, never as a pipeline component that could bias decisions.

\subsection{Why Human Groups Fail, and What That Implies}

The strongest reason to force a debate comes from old research on human groups. Janis studied policy fiascoes and named groupthink: cohesive groups suppress dissent, miss risks, and converge on confident wrong answers \cite{janis1972groupthink}. Tversky and Kahneman documented systematic judgment errors under uncertainty, including overconfidence \cite{tversky1974judgment}. Kahneman and Tversky's prospect theory showed that losses loom larger than equivalent gains, which explains why real traders hold losers and sell winners \cite{kahneman1979prospect}.

An LLM desk inherits these failure modes by analogy: agents prompted similarly tend to agree, and agreement feels like confidence. \sys{} counters structurally. The Bear Advocate exists to institutionalize dissent, and the Risk Manager exists to institutionalize loss aversion in the correct direction: it treats a bad trade as worse than a missed trade. The project also tests the failure directly: if Bull and Bear confidence always move together, the debate is decorative, and the hypothesis set measures exactly that.

% =====================================================================
\section{Risk: The Part That Says No}
\label{sec:risk}
% =====================================================================

The risk literature in this review is deliberately small and practical. The Sharpe ratio defines the reward-per-unit-of-risk number used to compare the simulated position series against baselines \cite{sharpe1994sharpe}. Maximum drawdown, the worst peak-to-trough loss, is the number the Risk Manager exists to limit, and its measurement discipline follows Lopez de Prado \cite{lopezdeprado2018advances}. The Deflated Sharpe Ratio \cite{bailey2014deflated} keeps the project honest about how good its risk-adjusted numbers really are.

Position sizing in \sys{} is capped, never optimized. The Risk Manager sets a maximum fraction of portfolio, reads Kronos's uncertainty width as its volatility signal, and can veto outright. The Pipeline then enforces the cap in code. This design inherits prospect theory's asymmetry as an explicit instruction: false positives, taking bad trades, cost more than false negatives, missing good ones \cite{kahneman1979prospect}.

\emph{What the project takes from this:} the risk layer is measured, not vibes. The evaluation checks whether the veto fires more in the stress window, whether caps tighten as uncertainty widens, and whether drawdown drops when the veto is active, using the no-veto ablation as the control.

% =====================================================================
\section{What the Literature Does Not Cover}
\label{sec:gap}
% =====================================================================

Reading across all of these areas, three things are missing, and together they are the space \sys{} occupies.

First, no published trading-agent system feeds a \emph{forecast distribution} from a dedicated candle foundation model into a risk layer. TradingAgents reads indicators with an LLM \cite{xiao2024tradingagents}; FinMem has no external forecasting component \cite{yu2023finmem}; Kronos forecasts but has no trading layer around it \cite{shi2025kronos}. The pieces exist side by side in the literature but have not been connected, and the connection is where the interesting question lives: what happens when a numeric uncertainty signal is allowed to overrule a language model's confident narrative?

Second, the debate literature measures debate on general reasoning tasks, not on trading decisions with money-shaped consequences. Whether an adversarial debate improves \emph{calibration} of confidence on trading calls, not just accuracy, is an open question. The Brier score and reliability analysis in the evaluation plan are designed to answer it.

Third, systems are evaluated by returns, not by auditability. For a decision-support product, the trail is the product. No paper in this review evaluates an agent pipeline on trail completeness, schema validity, and reproducibility as first-class outcomes. \sys{} does, because its contract is that a human reads and judges the reasoning before acting.

% =====================================================================
\section{How Each Paper Feeds the Project}
\label{sec:mapping}
% =====================================================================

Table~\ref{tab:mapping} compresses the review into a design map. Every row is a component of \sys{}, the papers it comes from, and the specific decision those papers drove.

\begin{table}[!t]
\centering
\caption{From literature to design: what each group of papers contributes to \sys{}.}
\label{tab:mapping}
\small
\begin{tabularx}{\columnwidth}{L{2.0cm} L{2.3cm} X}
\toprule
\textbf{Component} & \textbf{Key papers} & \textbf{What the project takes} \\
\midrule
Evaluation rules & Fama \cite{fama1970efficient}, Lo \cite{lo2004adaptive} & Buy-and-hold is the bar; results are window-specific; no profit claims. \\
\addlinespace[2pt]
Technical leg & Shi et al. (Kronos) \cite{shi2025kronos}, with Vaswani \cite{vaswani2017attention}, PatchTST \cite{nie2023patchtst}, GPT4TS \cite{zhou2023onefitsall}, Lag-Llama \cite{rasul2023lagllama}, Chronos \cite{ansari2024chronos} as lineage & Forecast raw candles with a foundation model; use the sampled spread, never invented chart narratives. \\
\addlinespace[2pt]
Sentiment leg & Lopez-Lira \cite{lopezlira2023chatgpt}, Tetlock \cite{tetlock2007giving}, Antweiler \cite{antweiler2004text}, Loughran \cite{loughran2011liability} & LLM reads headlines instead of dictionaries; works only from provided text; low confidence when news is thin. \\
\addlinespace[2pt]
Backtest design & Lopez de Prado \cite{lopezdeprado2018advances}, Bailey et al. \cite{bailey2014backtest,bailey2014deflated} & Walk-forward, no-lookahead loader, pre-registered ablations, deflated view of any Sharpe result. \\
\addlinespace[2pt]
Metrics & Sharpe \cite{sharpe1994sharpe}, MASE \cite{hyndman2006another}, Diebold-Mariano \cite{diebold1995comparing} & Every reported number has a source and a comparison target. \\
\addlinespace[2pt]
Agent roles & TradingAgents \cite{xiao2024tradingagents}, FinMem \cite{yu2023finmem}, surveys \cite{guo2024llmmultiagents,ding2024llmtrading}, ReAct \cite{yao2022react}, ChatDev \cite{qian2023chatdev}, FinRobot \cite{yang2024finrobot} & Role separation with fixed narrow jobs; no free-roaming autonomy; desk topology borrowed and then hardened. \\
\addlinespace[2pt]
Debate & Du \cite{du2023debate}, Liang \cite{liang2023divergent}, ChatEval \cite{chan2023chateval}, Smit \cite{smit2024mad} & Minimal two-turn adversarial format; rebuttal must name the Bull's claims; round count is a tunable, not a virtue. \\
\addlinespace[2pt]
Anti-groupthink & Janis \cite{janis1972groupthink}, Tversky \cite{tversky1974judgment}, Kahneman \cite{kahneman1979prospect} & Bear Advocate institutionalizes dissent; H4 measures whether disagreement actually exists. \\
\addlinespace[2pt]
Risk layer & Sharpe \cite{sharpe1994sharpe}, prospect theory \cite{kahneman1979prospect}, deflated Sharpe \cite{bailey2014deflated} & Veto-first design; uncertainty width drives sizing; cap enforced in code; bad trade weighted worse than missed trade. \\
\addlinespace[2pt]
Trust and checking & Hallucination survey \cite{ji2023hallucination}, LLM-as-judge \cite{zheng2023judging} & Schema-validated outputs, traceable numbers, judge used only in evaluation and checked against human labels. \\
\addlinespace[2pt]
Financial LLM context & BloombergGPT \cite{wu2023bloomberggpt}, FinGPT \cite{yang2023fingpt}, MarketSenseAI \cite{fatouros2024marketsenseai}, FinRL \cite{liu2023finrl} & LLMs read and combine evidence; code and specialized models compute; RL trading noted as the contrasting paradigm. \\
\addlinespace[2pt]
Orchestration & LangGraph \cite{langgraph2024} & Stateful graph with parallel fan-out and conditional veto edge. \\
\bottomrule
\end{tabularx}
\end{table}

% =====================================================================
\section{Summary}
% =====================================================================

The literature gives \sys{} everything it needs except the final connection. Finance says test honestly against simple baselines and expect little \cite{fama1970efficient,lopezdeprado2018advances}. Forecasting research contributes a candle foundation model whose sampled spread is a usable uncertainty signal \cite{shi2025kronos}. LLM research contributes evidence that models read financial text with real predictive value \cite{lopezlira2023chatgpt} and that role-structured teams behave better than one model doing everything \cite{xiao2024tradingagents,guo2024llmmultiagents}. Debate research says structured disagreement helps reasoning and warns that format and cost matter \cite{du2023debate,smit2024mad}. Psychology explains why a desk without a guaranteed dissenter drifts into groupthink \cite{janis1972groupthink}.

What no paper does is wire these together with an independent risk authority that reads forecast uncertainty and can veto the trade. That wiring, and the honest measurement of whether it improves calibration and drawdown, is the research content of \sys{}.

\bibliographystyle{IEEEtran}
\bibliography{references}

\end{document}
